% ===================================================================
%  तक्ता क्रमांक १२ — स्थानक. --- लोहमार्ग. --- नौका.
%
%  Traduction, faite en 2026, en marathi d'aujourd'hui, sur l'ido de
%  texto/io. Regles de la colonne : voir 10-takta-01.tex. Une seule
%  suite d'alineas numerotes : les cles ne portent pas d'indice de
%  scene. Le « Noto. » d'ouverture n'a pas de cle propre dans le
%  fac-simile : il fait partie du bloc d'apparat, et la traduction
%  le laisse la.
%
%  LE CHEMIN DE FER EST ARRIVE EN INDE AVANT D'ARRIVER DANS CE
%  LIVRET, ET LE MARATHI L'A NOMME. La premiere ligne indienne est
%  partie de Bori Bunder en 1853, soixante-treize ans avant le
%  Delmas-Tabeli, et elle partait justement du pays marathe. Le
%  vocabulaire n'a donc rien d'un calque :
%
%    फलाट (53)   le quai         रूळ (64)     les rails
%    डबा (37)    la voiture      मालगाडी (80) le train de marchandises
%    मालधक्का (71) la gare de marchandises
%    हमाल (16)   le porteur      वजनकाटा (21) la bascule
%    तिकीट (9)   le billet       खिडकी (8)    le guichet
%    वेळापत्रक (26) l'horaire     शिट्टी (59)  le sifflet
%    फाटा (78)   l'aiguillage    काटेवाला     l'aiguilleur
%    तार (5)     le telegramme   कारकून (22)  le commis
%
%  ET « तार » DIT LE TELEGRAMME PAR SON FIL, comme le francais dit
%  « le fil » et l'anglais « the wire ». Le mot marathe est celui du
%  fil de cuivre du tableau 5 : la meme chose a servi a nommer le
%  message qu'elle portait.
%
%  LE RESTE S'EMPRUNTE, ET C'EST LA MEME REGLE QUE PARTOUT : le
%  piston, le boiler, le tunnel, l'express n'ont pas de nom marathe
%  ancien parce que la machine a vapeur n'est pas d'ici. On ne
%  fabrique pas de mot pour boucher un trou.
%
%  DEUX INVERSIONS PREVENUES — le commis de la bascule et le pont de
%  fer sur le fleuve —, RESOLUES PAR UNE RELATIVE ET NON PAR UN
%  TIRET. Les deux remedes marchent, mais ils ne se valent pas :
%  quand le second terme qualifie durablement le premier (le commis
%  QUI TIENT la bascule, le pont QUI FRANCHIT le fleuve), la
%  relative se lit comme du marathi ordinaire ; le tiret, lui, sert
%  quand le second terme n'est qu'une precision de lieu ou de
%  compagnie, jetee apres coup. Trois autres inversions ont passe
%  quand meme, et renvoji.py les a vues.
% ===================================================================

%%K t12-apar-1 apar
\VUsaut{4.0mm}
\VUtitre{15.0pt}{60}{तक्ता क्रमांक १२}
\VUsaut{2.4mm}
\VUpk{11.4pt}{40}{स्थानक. --- लोहमार्ग. --- नौका.}
\VUsaut{3.4mm}
\textsc{टीप.} --- \textit{प्रत्येक देशात वापरली जाणारी वेळापत्रके वेगवेगळी असतात, म्हणून आपण उदाहरणादाखल} \VUgras{फ्रेंच वेळापत्रकाच्या} \textit{वेळा वापरतो.}

%%K t12-01-1 p
१. --- मी नुकताच जर्मनीतून प्रवास केला. निघण्यापूर्वी गाड्यांच्या सुटण्याच्या वेळा कळाव्यात
म्हणून मी \VUgras{वेळापत्रकाचे पुस्तक} \textsuperscript{(15)} विकत घेतले. सर्वात सोयीची गाडी
पॅरिसहून रात्री नऊ वाजता सुटते आणि स्ट्रासबूर्गला एक वाजून तेरा मिनिटांनी पोहोचते हे पाहून मी
माझ्या प्रवासाची तयारी केली, आणि दुसऱ्या दिवशी मी स्थानकाकडे निघालो.

%%K t12-02-1 p
२. --- मी प्रस्थानाच्या मोठ्या दालनात शिरलो तेव्हा एका म्हाताऱ्या गावच्या
\VUgras{धर्मगुरूमागोमाग} \textsuperscript{(2)} — ज्याने हातात आपली \VUgras{प्रवासाची पिशवी}
\textsuperscript{(3)} घेतली होती — पुष्कळ \VUgras{प्रवासी} \textsuperscript{(1)} आधीच पोहोचले
होते.

%%K t12-02-2 p
मला एक \VUgras{तार} \textsuperscript{(5)} पाठवायची होती, म्हणून एका \VUgras{तपासनिसाकडून}
\textsuperscript{(6)} मी \VUgras{तार कार्यालय} \textsuperscript{(4)} विचारून घेतले.

%%K t12-03-1 p
३. --- \VUgras{प्रतीक्षालयात} \textsuperscript{(7)} जाण्यापूर्वी मी परतीच्या प्रवासासह
\VUgras{तिकीट} \textsuperscript{(9)} काढून घेतले.

%%K t12-04-1 p
४. --- \VUgras{खिडकीशी} \textsuperscript{(8)} मला फार वेळ थांबावे लागले नाही, कारण तिथे थोडीच
माणसे होती. रजेवर निघालेल्या एका \VUgras{शिपायाने} \textsuperscript{(10)} सौजन्याने मला आपली
पाळी दिली, त्यामुळे \VUgras{स्थानकप्रमुखांकडून} \textsuperscript{(13)} थोडी माहिती विचारायला
मला पुरेसा वेळ मिळाला — ते देखरेखीच्या \VUgras{आयुक्तांशी} \textsuperscript{(12)} आणि
कर्तव्यावरच्या \VUgras{पोलिसाशी} \textsuperscript{(11)} बोलत होते.

%%K t12-tit-1 sub
\VUpk{10.2pt}{40}{सामानाची नोंदणी.}
\VUsaut{3.0mm}

%%K t12-05-1 p
५. --- \VUgras{पुस्तकांच्या गाळ्यावर} \textsuperscript{(14)} वर्तमानपत्र विकत घेऊन मी माझे
\VUgras{सामान} \textsuperscript{(17)} वजन करून घेतले, जे \VUgras{हमालाने}
\textsuperscript{(16)} \VUgras{सामानाच्या गाडीवरून} \textsuperscript{(19)} नुकतेच आणले होते;
\VUgras{त्या कारकुनाने} \textsuperscript{(22)}, जो \VUgras{वजनकाट्याचा} \textsuperscript{(21)}
कारभार पाहतो, मला सांगितले की माझी ट्रंक पस्तीस किलोग्रॅम भरते; म्हणजे पाच किलोग्रॅम
\VUgras{जास्तीचे वजन} झाले, ज्यासाठी मला \VUgras{सामानाच्या खिडकीवर} \textsuperscript{(23)}
जादा पैसे भरावे लागतील.

%%K t12-06-1 p
६. --- मी फार लवकर आलो असल्यामुळे स्थानकातली प्रवाशांची वर्दळ न्याहाळायला मला मोकळा वेळ
मिळाला. काही जण आपले किरकोळ सामान \VUgras{सामान ठेवण्याच्या खोलीत} \textsuperscript{(25)} ठेवत
किंवा परत घेत होते; दुसरे \VUgras{वेळापत्रके} \textsuperscript{(26)} पाहत होते. येणारे प्रवासी
\VUgras{बाहेर पडण्याच्या मार्गाकडे} \textsuperscript{(32)} घाई करत होते, हातात \VUgras{सुटकेस}
\textsuperscript{(24)} घेऊन. काही जण \VUgras{सामान वाटपाच्या खोलीत} \textsuperscript{(27)}
थांबले होते आणि आपली \VUgras{पावती} \textsuperscript{(29)} दाखवत होते, आपल्या ट्रंका,
\VUgras{पेट्या} \textsuperscript{(28)} किंवा \VUgras{टोपल्या} \textsuperscript{(30)} परत
घ्यायला.

%%K t12-07-1 p
७. --- \VUgras{जकात अधिकारी} \textsuperscript{(33)} काही जाहीर करण्यासारखे आहे का हे पाहायला
गाठोडी काळजीपूर्वक तपासत होते.

%%K t12-08-1 p
८. --- काही प्रवासी \VUgras{उपाहारगृहात} \textsuperscript{(34)} गेले, जिथे एक \VUgras{वाढपी}
\textsuperscript{(35)} त्यांना वाढायची लगबग करत होता. त्यांना घाई करावी लागते, कारण ती
\VUgras{गाडी} \textsuperscript{(36)} जलद गाडी असल्यामुळे स्थानकात फार वेळ थांबणार नाही.

%%K t12-09-1 p
९. --- मग मी प्रस्थानाच्या फलाटावर गेलो आणि प्रवासात डबा बदलावा लागू नये म्हणून थेट जाणारा
\VUgras{डबा} \textsuperscript{(37)} शोधला. मी एका बोळाच्या डब्यात चढलो, ज्यात धूम्रपान
करणाऱ्यांसाठी एक \VUgras{कप्पा} \textsuperscript{(38)} आणि « एकट्या बायकांसाठी » राखीव एक
कप्पा होता.

%%K t12-tit-2 sub
\VUpk{10.2pt}{40}{रात्रीचे प्रस्थान.}
\VUsaut{3.0mm}

%%K t12-10-1 p
१०. --- \VUgras{पायरी} \textsuperscript{(40)} चढून, \VUgras{दार} \textsuperscript{(39)} लावून,
\VUgras{पडदा} \textsuperscript{(41)} गुंडाळून आणि माझे किरकोळ सामान \VUgras{जाळीवर}
\textsuperscript{(43)} ठेवून मी \VUgras{बाकड्यावर} \textsuperscript{(42)} बसलो.
\VUgras{दिवेवाल्याला} \textsuperscript{(44)} दिवे लावताना पाहून मला वाटले की मला रात्र डब्यातच
काढावी लागणार, आणि \VUgras{उश्यांच्या} \textsuperscript{(45)} भाड्याचे काम पाहणाऱ्या
कर्मचाऱ्याला बोलावून मी एक मागितली.

%%K t12-11-1 p
११. --- तिकीट तपासनीस तिकिटे तपासायला आला. वेळ बरोबर झाली असूनही आपण अजून का निघालो नाही असे
मी त्याला विचारले. त्याने मला सांगितले की \VUgras{गाडीप्रमुख} \textsuperscript{(46)}
\VUgras{सामानाच्या डब्यात} \textsuperscript{(47)} सायकली ठेवण्यात गुंतले आहेत. शिवाय, अजून
सगळे पाय शेकायचे डबे \textsuperscript{(48)} बदलले गेलेले नाहीत.

%%K t12-12-1 p
१२. --- पण आम्ही लवकरच निघणार होतो, कारण \VUgras{कोळसा घालणाऱ्याने} \textsuperscript{(50)}
आपल्या \VUgras{कोळशाच्या डब्यावर} \textsuperscript{(49)} \VUgras{इंजिनाची}
\textsuperscript{(54)} आग फुलवायला कोळसा उचलला होता, आणि \VUgras{इंजिनचालक}
\textsuperscript{(51)} फक्त \VUgras{स्थानकाच्या उपप्रमुखांच्या} \textsuperscript{(52)}
इशाऱ्याची वाट पाहत होता, जे \VUgras{फलाटावर} \textsuperscript{(53)} होते.

%%K t12-13-1 p
१३. --- इंजिनाचा \VUgras{बॉयलर} \textsuperscript{(56)} मंद गुरगुरत होता; \VUgras{धुराड्यातून}
\textsuperscript{(55)} \VUgras{वाफेत} \textsuperscript{(60)} मिसळलेल्या धुराचे लोट बाहेर पडत
होते. कर्कश \VUgras{शिट्टी} \textsuperscript{(59)} घुमली आणि \VUgras{दट्टे}
\textsuperscript{(58)} हलू लागले, त्या जड यंत्राची \VUgras{चाके} \textsuperscript{(57)} ढकलत.

%%K t12-tit-3 sub
\VUsaut{3.0mm}
\VUpk{10.2pt}{40}{निघाली!}
\VUsaut{3.0mm}

%%K t12-14-1 p
१४. --- \VUgras{रुळांवरून} \textsuperscript{(64)} धावणाऱ्या गाडीचा वेग वाढत गेला;
\VUgras{फलाट} \textsuperscript{(65)} नाहीसे झाले; आम्ही \VUgras{स्वच्छतागृहे}
\textsuperscript{(66)} मागे टाकली, आणि दुसऱ्या \VUgras{मार्गावर} \textsuperscript{(63)} उभी
केलेली डब्यांची रांगही : \VUgras{झोपण्याचे डबे} \textsuperscript{(67)}, \VUgras{भोजनाचा डबा}
\textsuperscript{(68)}, \VUgras{टपालाचा डबा} \textsuperscript{(69)}.

%%K t12-noto-1 noto
\VUnotes{143.2mm}{%
\textit{(*) टीप.} --- \textit{अर्थातच प्रवासाचे वर्णन युद्धापूर्वीच्या सरहद्दीप्रमाणे आहे.}%
}

%%K t12-15-1 p
१५. --- मग \VUgras{मालधक्का} \textsuperscript{(71)} आमच्या डोळ्यांसमोरून सरकला. छपरांखाली
कर्मचारी संथ गाडीने पाठवलेला \VUgras{माल} \textsuperscript{(72)} भरत होते.
\VUgras{डबे जोडणारे} \textsuperscript{(76)} \VUgras{पाण्याच्या टाकीजवळ} \textsuperscript{(73)}
\VUgras{जनावरांचे डबे} \textsuperscript{(75)} आणि \VUgras{सपाट डबे} \textsuperscript{(79)}
हलवत होते, \VUgras{मालगाडी} \textsuperscript{(80)} तयार करण्यासाठी.

%%K t12-16-1 p
१६. --- आम्ही शेवटचा \VUgras{फाटा} \textsuperscript{(78)} ओलांडला, ज्याच्याजवळ काटेवाला उभा
होता; आम्ही \VUgras{संकेतचकती} \textsuperscript{(86)} मागे टाकली आणि स्थानक दूर नाहीसे झाले.

%%K t12-17-1 p
१७. --- लवकरच आम्ही तो \VUgras{लोखंडी पूल} \textsuperscript{(74)} ओलांडला, जो \VUgras{नदीवरून}
\textsuperscript{(81)} जातो. तो पूल ओलांडल्यावर आम्ही नदीकाठाने गेलो, जिच्यावर जबरदस्त
\VUgras{ओढनौका} \textsuperscript{(82)} जड \VUgras{मालवाहू नौका} \textsuperscript{(83)} ओढत
होत्या. शिवाय आम्हाला एक \VUgras{प्रवासी नौकाही} \textsuperscript{(84)} दिसली, जेव्हा ती एका
\VUgras{तराफ्याला} \textsuperscript{(85)} लागत होती.

%%K t12-18-1 p
१८. --- कित्येक स्थानके भरधाव मागे टाकून आम्ही काही \VUgras{बोगदे} \textsuperscript{(87)}
ओलांडले आणि \VUgras{फाटकवाल्यांची} असंख्य \VUgras{घरे} \textsuperscript{(88)} मागे सोडली.
गाडीच्या हेलकाव्यांनी झोका दिल्यासारखे होऊन मी गाढ झोपी गेलो.

%%K t12-tit-4 sub
\VUsaut{3.0mm}
\VUpk{10.2pt}{40}{डॉइच-आव्रिकूर स्थानक.}
\VUsaut{3.0mm}

%%K t12-19-1 p
१९. --- « डॉइच-आव्रिकूर! \textsuperscript{(*)}. सगळे उतरा! » असे ओरडणाऱ्या एका कर्मचाऱ्याच्या
आवाजाने मी एकदम जागा झालो. जकातीची तपासणी झाली आणि सरहद्दीचे सगळे कंटाळवाणे सोपस्कारही. मला
माझे खिशातले घड्याळ स्थानकाच्या \VUgras{मोठ्या घड्याळाशी} \textsuperscript{(89)} जुळवावे
लागले, जे पंचावन्न मिनिटे पुढे होते; नुकताच जागा झाल्यामुळे मला \VUgras{मोठा काटा}
\textsuperscript{(90)} \VUgras{लहानापासून} \textsuperscript{(91)} कष्टानेच ओळखता आला; सकाळच्या
धुक्यामुळे \VUgras{तबकडीच} \textsuperscript{(92)} पार अस्पष्ट झाली होती.

%%K t12-20-1 p
२०. --- जकात अधिकारी तपासणी करत असताना मी जांभया देत स्थानकाच्या भिंती सजवणाऱ्या
\VUgras{जाहिराती} \textsuperscript{(93)} वाचल्या. वेळ घालवायला मी न्याहारी केली, स्ट्रासबूर्ग
अजून किती दूर आहे हे पाहायला जर्मन \VUgras{जाळ्याचा} \textsuperscript{(94)} नकाशा पाहिला, आणि
माझ्या प्रवासाचे ध्येय लवकरच गाठू या आशेने बळ मिळून मी माझ्या डब्यात परत शिरलो.

\VUsaut{6.0mm}
\VUfilet{22mm}
